\subsection{文件读写与删除}
在src文件夹中的file.c文件中实现文件的读写与删除功能,具体包括文件的创建、读写和删除操作,具体代码分别如下:
\begin{itemize}
    \item 文件创建函数:
    \begin{ccode}{文件创建}
// 创建文件，写文件
bool fs_create(const char*filename){
    for (int i = 0; i < MAX_ENTRY; i++){
        if(Directory[i].ac == 0){
            // 找到空闲块
            strncpy(Directory[i].name,filename,DIR_NAME_LENGTH);
            Directory[i].size = 0;
            Directory[i].start_block = -1;
            Directory[i].ac = 1; 
            Directory[i].type = 0;//0表示文件
            Directory[i].parent = current_dir; //设置父目录索引
            printf("File created successfully : %s\n",filename);
            return true;
        }
    }
    printf("Failed to created file : %s\n",filename);
    return false;
}
    \end{ccode}
    文件创建就是接收参数文件名filename,在Directory数组中找一个空闲块,并设置该块对应属性,包括文件名、文件大小、文件类型、是否被分配以及父目录索引。
    \item 写文件函数:
    \begin{ccode}{写文件函数}
int fs_write(int fd,const char* data,int size){
    if(fd < 0 || fd >= MAX_FD || data == NULL || size < 0) return -1;
    if(FDTable[fd].used == 0){
        printf("write failed : fd is not open\n");
        return -1;
    }

    if(FDTable[fd].mode != MODE_WRITE && FDTable[fd].mode != MODE_RDWR && FDTable[fd].mode != MODE_WAPPEND){
        printf("write failed : fd is not writable");
        return -1;
    }

    if (size == 0){
        return 0;
    }

    int dir_index = FDTable[fd].dir_index;
    int offset = FDTable[fd].offset;
    int end_pos = offset + size;

    int need_blocks = (end_pos + BLOCK_SIZE - 1) / BLOCK_SIZE;

    // 确保文件有足够多的 block
    if (Directory[dir_index].start_block == -1) {
        int new_block = -1;

        for (int i = 0; i < BLOCK_NUM; i++) {
            if (FAT[i] == 0) {
                new_block = i;
                break;
            }
        }

        if (new_block == -1) {
            printf("write failed: no free block\n");
            return -1;
        }

        FAT[new_block] = -1;
        memset(Disk + new_block * BLOCK_SIZE, 0, BLOCK_SIZE);
        Directory[dir_index].start_block = new_block;
    }

    int cur = Directory[dir_index].start_block;
    int block_count = 1;

    while (FAT[cur] != -1) {
        cur = FAT[cur];
        block_count++;
    }

    while (block_count < need_blocks) {
        int new_block = -1;

        for (int i = 0; i < BLOCK_NUM; i++) {
            if (FAT[i] == 0) {
                new_block = i;
                break;
            }
        }

        if (new_block == -1) {
            printf("write failed: no enough space\n");
            return -1;
        }

        FAT[new_block] = -1;
        memset(Disk + new_block * BLOCK_SIZE, 0, BLOCK_SIZE);

        FAT[cur] = new_block;
        cur = new_block;
        block_count++;
    }

    //找offset对应的block
    int block_index = offset / BLOCK_SIZE;
    int block_offset = offset % BLOCK_SIZE;

    cur = Directory[dir_index].start_block;

    for(int i = 0; i < block_index;i++){
        cur = FAT[cur];
    }

    // 从offset 开始写数据
    int bytes_written = 0;
    while (cur != -1 && bytes_written < size){
        int write_size = BLOCK_SIZE - block_offset;

        if(write_size > size - bytes_written){
            write_size = size - bytes_written;
        }
        memcpy(
            Disk + cur * BLOCK_SIZE + block_offset,
            data + bytes_written,
            write_size
        );
        bytes_written += write_size;
        block_offset = 0;
        if(bytes_written < size){
            cur = FAT[cur];
        }
    }

    //更新offset和文件大小

    FDTable[fd].offset += bytes_written;
    if(FDTable[fd].offset > Directory[dir_index].size){
        Directory[dir_index].size = FDTable[fd].offset;
    }

    return bytes_written;
}
    \end{ccode}
    写文件函数接收参数文件描述符fd,数据指针data以及数据大小size,实际写入前需要确保fd有效且文件可写,同时确保FAT中有足够的多的块写入内容(否则写文件失败),最后根据offset确认写入的块的位置并写入data。
    \item 读文件函数:
    \begin{ccode}{读文件函数}
int fs_read(int fd,char *buffer,int size){

    if (fd < 0 || fd >= MAX_FD || buffer == NULL || size < 0) {
        return -1;
    }

    if (FDTable[fd].used == 0) {
        printf("read failed: file is not open\n");
        return -1;
    }

    if (FDTable[fd].mode != MODE_READ && FDTable[fd].mode != MODE_RDWR ){
        printf("read failed : file is not readable\n");
        return -1;
    }
    int dir_index = FDTable[fd].dir_index;
    int file_size = Directory[dir_index].size;
    int offset = FDTable[fd].offset;

    if (offset >= file_size) {
        buffer[0] = '\0';
        return 0;
    }

    int bytes_to_read = size;
    if (offset + bytes_to_read > file_size) {
        bytes_to_read = file_size - offset;
    }

    int bytes_read = 0;
    int cur = Directory[dir_index].start_block;
    int skip_blocks = offset / BLOCK_SIZE;
    int block_offset = offset % BLOCK_SIZE;

    // 跳到 offset 所在的块
    for (int i = 0; i < skip_blocks && cur != -1; i++) {
        cur = FAT[cur];
    }

    while (cur != -1 && bytes_read < bytes_to_read) {
        int read_size = BLOCK_SIZE - block_offset;

        if (read_size > bytes_to_read - bytes_read) {
            read_size = bytes_to_read - bytes_read;
        }

        memcpy(
            buffer + bytes_read,
            Disk + cur * BLOCK_SIZE + block_offset,
            read_size
        );

        bytes_read += read_size;
        block_offset = 0;
        cur = FAT[cur];
    }

    buffer[bytes_read] = '\0';
    FDTable[fd].offset += bytes_read;

    return bytes_read;
}
    \end{ccode}
    读文件函数接收参数文件描述符fd,数据指针buffer以及要读的数据大小size,类似写文件函数,同样要确保fd有效且文件可读,同时根据offset确认读取的块的位置并读取数据。
    \item 删除文件函数:
    \begin{ccode}{删除文件函数}
bool fs_delete(const char *filename){
    for (int i = 0 ; i < MAX_ENTRY; i ++){
        if(Directory[i].ac == 1 && Directory[i].parent == current_dir && Directory[i].type == 0 && strcmp(Directory[i].name,filename) == 0){
            //文件占用的磁盘块在FAT中全部标记为0
            int cur = Directory[i].start_block;
            while (cur != -1){
                int next = FAT[cur];
                FAT[cur] = 0;
                if(next == -1){//找到末尾块
                    break;
                }
                cur = next;
            }
            Directory[i].ac = 0;
            printf("File deleted successfully : %s \n",filename);
            return true;
        }
    }

    printf("Failed to delete file: %s\n", filename);
    return false;
}
    \end{ccode}
    删除文件函数类似文件创建函数,接收参数文件民名filename,然后在Directory数组中查找该文件,找到后释放该文件占用的所有FAT块(从start\_block开始全部置0),并设置为未分配状态即可。
   \end{itemize}
